\section{MỞ ĐẦU VÀ TỔNG QUAN ĐỀ TÀI}\label{sec:mo_dau_tong_quan}

\subsection{Lý do chọn đề tài}

Mạng Campus là hạ tầng nền tảng của trường đại học, phục vụ học tập, nghiên cứu, quản trị và các dịch vụ số. Khi trường mở thêm cơ sở tại nhiều địa phương, hệ thống mạng phải đồng thời bảo đảm hiệu năng, tính sẵn sàng, bảo mật, khả năng quản trị và mở rộng. Mạng truyền thống được cấu hình riêng trên từng thiết bị: mỗi switch tự học địa chỉ và tự chạy Spanning Tree Protocol (STP), kết nối giữa các cơ sở phụ thuộc một đường truyền thuê riêng trên mạng diện rộng Wide Area Network (WAN). Mỗi thay đổi Virtual Local Area Network (VLAN) hay chính sách phải làm trên nhiều thiết bị, dễ sai lệch, thiếu cái nhìn toàn cục và không tự chọn đường theo chất lượng.

Hai công nghệ Software-Defined Networking (SDN) và Software-Defined Wide Area Network (SD-WAN) giải quyết trực tiếp các hạn chế này. SDN tách mặt phẳng điều khiển khỏi thiết bị chuyển mạch để một Controller lập trình, giám sát và áp chính sách tập trung cho mạng nội bộ. SD-WAN tạo mạng overlay mã hóa trên nhiều đường truyền như Internet và Multiprotocol Label Switching (MPLS), quản lý tập trung và chọn đường theo chất lượng thực tế. Đó là lý do nhóm chọn đề tài ``Xây dựng Campus Network sử dụng SDN và SD-WAN''.

\subsection{Mục tiêu thực hiện đề tài}

Đề tài thiết kế, triển khai và kiểm chứng một mô hình mạng Campus đa cơ sở, trong đó mạng nội bộ cơ sở chính do SDN điều khiển và các cơ sở được kết nối bằng SD-WAN. Các mục tiêu cụ thể:

\begin{itemize}
    \item Thiết kế Campus chính theo mô hình Core--Distribution--Access với gateway dự phòng Virtual Router Redundancy Protocol (VRRP), định tuyến Open Shortest Path First (OSPF), cặp tường lửa Cisco Adaptive Security Virtual Appliance (ASAv) chạy High Availability (HA), vùng Demilitarized Zone (DMZ) và Server Farm; thiết kế mô hình chi nhánh rút gọn Firewall-as-Core cho ba cơ sở Cần Thơ, Đà Nẵng, Nha Trang.
    \item Xây dựng quy hoạch VLAN, địa chỉ IP, System-IP, Transport Locator (TLOC) và số hiệu hệ tự trị Autonomous System Number (ASN) thống nhất cho toàn hệ thống.
    \item Phát triển ứng dụng Controller Ryu theo OpenFlow 1.3 điều khiển toàn bộ switch lớp 2 Open vSwitch (OVS) tại cơ sở chính: học MAC theo VLAN, cây chuyển tiếp tính tập trung, tự phục hồi, kênh quản trị VLAN 99 không phụ thuộc STP, thêm VLAN và áp chính sách qua Application Programming Interface (API), cùng giao diện giám sát Campus SDN Console.
    \item Triển khai Cisco SD-WAN (Viptela) kết nối các site qua hai nhà cung cấp Internet và MPLS với underlay Border Gateway Protocol (BGP), xác thực bằng Public Key Infrastructure (PKI), giao thức Overlay Management Protocol (OMP), giám sát đường hầm bằng Bidirectional Forwarding Detection (BFD), chính sách Service Level Agreement (SLA), định tuyến theo ứng dụng Application-Aware Routing (AAR) và cảnh báo sự cố qua email.
    \item Cho người dùng các cơ sở ra Internet tại chỗ theo Direct Internet Access (DIA) qua Network Address Translation (NAT), có dự phòng; cung cấp các dịch vụ dùng chung Dynamic Host Configuration Protocol (DHCP), Domain Name System (DNS), Web, Mail, Syslog và bảo vệ bằng Access Control List (ACL) nhiều lớp.
    \item Kiểm thử theo lớp, đo các tiêu chí đánh giá của đề cương và ghi nhận trung thực các hạn chế.
\end{itemize}

\subsection{Đối tượng và phạm vi nghiên cứu}

\textbf{Đối tượng nghiên cứu} là hệ thống mạng Campus đa cơ sở trong môi trường đại học: chuyển mạch và phân đoạn VLAN, dự phòng lớp 2 và lớp 3, định tuyến nội bộ và liên miền, tường lửa và DMZ, SDN với OpenFlow, SD-WAN và các dịch vụ mạng nền tảng.

\textbf{Phạm vi thực hiện:}
\begin{itemize}
    \item Mô hình được hiện thực trên nền tảng mô phỏng Emulated Virtual Environment -- Next Generation (EVE-NG) với 74 node: Campus chính (Site 100), ba cơ sở phụ (Site 200, 300, 400), trung tâm điều khiển SD-WAN (Site 900), hai router mô phỏng nhà cung cấp Internet, MPLS và một cơ sở mở rộng (Site 500). Switch Cisco dùng image IOS on Linux (IOL).
    \item SDN áp dụng cho tầng Distribution và Access của Campus chính (sáu OVS); tầng Core và tường lửa giữ cơ chế truyền thống theo mô hình Hybrid SDN. Liên kết giữa các cơ sở do SD-WAN đảm nhận.
    \item Mạng WAN mô phỏng dùng dải \texttt{203.0.113.0/24} và \texttt{100.64.0.0/10}; router nhà cung cấp Internet nối ra mạng thật của phòng lab. Mạng không dây và xác thực 802.1X nằm ngoài phạm vi.
    \item Đánh giá theo bảy nhóm tiêu chí của đề cương (hiệu năng, mở rộng, độ tin cậy, quản trị, bảo mật, Quality of Service (QoS), quy mô); thông lượng và độ trễ đo trong môi trường ảo hóa nên không phản ánh hiệu năng thiết bị thật.
\end{itemize}

\subsection{Phương pháp nghiên cứu}

\begin{itemize}
    \item \textbf{Nghiên cứu tài liệu:} các chuẩn và tài liệu kỹ thuật về OSPF, BGP, VRRP, BFD, OpenFlow 1.3, OVS, Ryu, Cisco ASA và Cisco SD-WAN.
    \item \textbf{Thiết kế:} phân tích yêu cầu, xây dựng kiến trúc, sơ đồ vật lý và sơ đồ logic, quy hoạch địa chỉ, định tuyến, SDN, SD-WAN, dịch vụ và bảo mật.
    \item \textbf{Hiện thực:} cấu hình từ dưới lên (underlay, controller, hạ tầng, dịch vụ, chính sách); phát triển ứng dụng Ryu bằng Python; quản lý cấu hình mọi thiết bị bằng Git.
    \item \textbf{Kiểm thử theo lớp:} liên kết, lớp 2, lớp 3, dịch vụ, ứng dụng; mọi kết luận dựa trên đầu ra lệnh, bộ đếm hoặc bắt gói; các phép đo được tự động hóa bằng script để có thể chạy lại.
\end{itemize}

\subsection{Ý nghĩa thực tiễn và bố cục báo cáo}

Mô hình cho thấy cách kết hợp SDN và SD-WAN với các cơ chế truyền thống trong một hệ thống thống nhất; các nguyên nhân gốc của sự cố đã gặp (phụ thuộc vòng giữa switch và Controller trong SDN in-band, quy tắc in-band ẩn gây bão broadcast, định tuyến tới controller cho từng màu TLOC, xung đột tên miền \texttt{.local} với mDNS) là kinh nghiệm có giá trị cho triển khai thực tế. Mô hình cũng có thể dùng làm nền tảng thực hành cho các học phần mạng máy tính.

Báo cáo gồm bốn chương:
\begin{itemize}
    \item \textbf{Chương 1 -- Mở đầu và tổng quan đề tài:} lý do, mục tiêu, phạm vi, phương pháp và bố cục.
    \item \textbf{Chương 2 -- Cơ sở lý thuyết:} mạng Campus phân cấp, VLAN, dự phòng, định tuyến, dịch vụ, tường lửa, SDN, SD-WAN và EVE-NG.
    \item \textbf{Chương 3 -- Mô hình mạng Campus đề xuất:} yêu cầu, kiến trúc, sơ đồ vật lý và sơ đồ logic, quy hoạch địa chỉ, thiết kế định tuyến, SDN, SD-WAN, dịch vụ và bảo mật.
    \item \textbf{Chương 4 -- Triển khai, đánh giá và kết luận:} hiện thực, các vấn đề kỹ thuật, kết quả kiểm thử, đánh giá theo tiêu chí của đề cương, kết luận và hướng phát triển.
\end{itemize}
